% GENERATED FILE. Edit canonical lessons, then run npm run generate:books.
% canonical-source-hash: fnv1a64:41671e36f7871466
% canonical-lessons: IT-C277-durare, IT-C277-durata, IT-C277-istante, IT-C277-attimo, IT-C277-fase

\chapter{Quanto è durato}
\label{ch:it-quanto-e-durato}
\hlchaptermodality{277}

\begin{chapteropening}
\textbf{By the end of this chapter:} \emph{I can say how long something lasted, and name a moment, an instant and a stage.}

Complete the last lesson of chapter 277: I can say how long something lasted, and name a moment, an instant and a stage.
\end{chapteropening}

\section[durare]{durare — to last}

\label{lesson:IT-C277-durare}



\begin{quote}
Before the new one: say the Italian for the chronology, then the Italian for the sequence.
\end{quote}

\subsection*{You'll want to know: durare}
\textbf{durare} — \textquotedblleft{}to last\textquotedblright{}.

\begin{quote}
Il film è durato tre ore.
\end{quote}

\begin{quote}
\emph{The film lasted three hours.}
\end{quote}

In the \emph{passato prossimo}, \emph{è durato} is the usual form; \emph{ha durato} is also heard.

\begin{cousinweb}
\emph{durare} is inherited from Latin \emph{dūrāre} \textquotedblleft{}to harden; to hold out, to last\textquotedblright{} (EtymDB), built on \emph{dūrus} \textquotedblleft{}hard\textquotedblright{} (Lewis \& Short).
\end{cousinweb}

\begin{grammarlens}[title={What you've built}]
One more word for telling, explaining and describing at B1.
\end{grammarlens}

\subsection*{Your turn}
\begin{itemize}
\raggedright
  \item \emph{Say it:} \emph{durare}
  \item \emph{Say it:} \emph{durare}, once more
  \item \emph{Say it:} a sentence of your own with \emph{durare}
  \item \emph{Recall:} say the Italian for the chronology, then the Italian for the sequence, then say \emph{durare} again
\end{itemize}

\begin{tcolorbox}[breakable,colback=teal!4,colframe=teal!35!black,title={Before you move on}]
What does durare mean? (\textquotedblleft{}To last\textquotedblright{}.) Say it once more.
\end{tcolorbox}
\section[la durata]{la durata — the length of time, the duration}

\label{lesson:IT-C277-durata}



\begin{quote}
Before the new one: say the Italian for the sequence, then the Italian for to last.
\end{quote}

\subsection*{You'll want to know: la durata}
\textbf{la durata} — \textquotedblleft{}the length of time, the duration\textquotedblright{}.

\begin{quote}
La durata del viaggio è di due ore.
\end{quote}

\begin{quote}
\emph{The journey lasts two hours.}
\end{quote}

Built on \emph{durare}, ‘to last’.

\begin{cousinweb}
\emph{durata} is built on \emph{durare} \textquotedblleft{}to last\textquotedblright{}, from Latin \emph{dūrāre} (EtymDB).
\end{cousinweb}

\begin{grammarlens}[title={What you've built}]
One more word for telling, explaining and describing at B1.
\end{grammarlens}

\subsection*{Your turn}
\begin{itemize}
\raggedright
  \item \emph{Say it:} \emph{la durata}
  \item \emph{Say it:} \emph{la durata}, once more
  \item \emph{Say it:} a sentence of your own with \emph{la durata}
  \item \emph{Recall:} say the Italian for the sequence, then the Italian for to last, then say \emph{la durata} again
\end{itemize}

\begin{tcolorbox}[breakable,colback=teal!4,colframe=teal!35!black,title={Before you move on}]
What does la durata mean? (\textquotedblleft{}The length of time, the duration\textquotedblright{}.) Say it once more.
\end{tcolorbox}
\section[l'istante]{l'istante — the instant}

\label{lesson:IT-C277-istante}



\begin{quote}
Before the new one: say the Italian for to last, then the Italian for the length of time.
\end{quote}

\subsection*{You'll want to know: l'istante}
\textbf{l'istante} — \textquotedblleft{}the instant\textquotedblright{}.

\begin{quote}
In un istante la stanza era vuota.
\end{quote}

\begin{quote}
\emph{In an instant the room was empty.}
\end{quote}

\emph{All'istante} means ‘at once, immediately’.

\begin{cousinweb}
\emph{istante} is a learned borrowing from Latin \emph{īnstāns} \textquotedblleft{}pressing, present\textquotedblright{}, from \emph{īnstāre} \textquotedblleft{}to stand upon, to press\textquotedblright{} (EtymDB): the moment that presses on you now.
\end{cousinweb}

\begin{grammarlens}[title={What you've built}]
One more word for telling, explaining and describing at B1.
\end{grammarlens}

\subsection*{Your turn}
\begin{itemize}
\raggedright
  \item \emph{Say it:} \emph{l'istante}
  \item \emph{Say it:} \emph{l'istante}, once more
  \item \emph{Say it:} a sentence of your own with \emph{l'istante}
  \item \emph{Recall:} say the Italian for to last, then the Italian for the length of time, then say \emph{l'istante} again
\end{itemize}

\begin{tcolorbox}[breakable,colback=teal!4,colframe=teal!35!black,title={Before you move on}]
What does l'istante mean? (\textquotedblleft{}The instant\textquotedblright{}.) Say it once more.
\end{tcolorbox}
\section[l'attimo]{l'attimo — the moment, a split second}

\label{lesson:IT-C277-attimo}



\begin{quote}
Before the new one: say the Italian for the length of time, then the Italian for the instant.
\end{quote}

\subsection*{You'll want to know: l'attimo}
\textbf{l'attimo} — \textquotedblleft{}the moment, a split second\textquotedblright{}.

\begin{quote}
Aspetta un attimo, arrivo!
\end{quote}

\begin{quote}
\emph{Wait a moment, I am coming!}
\end{quote}

\emph{Un attimo!} — ‘One moment!’ — is heard everywhere.

\begin{cousinweb}
\emph{attimo} is usually traced to Latin \emph{atomus}, from Greek \emph{átomos} \textquotedblleft{}that cannot be cut\textquotedblright{} (Lewis \& Short): the smallest slice of time.
\end{cousinweb}

\begin{grammarlens}[title={What you've built}]
One more word for telling, explaining and describing at B1.
\end{grammarlens}

\subsection*{Your turn}
\begin{itemize}
\raggedright
  \item \emph{Say it:} \emph{l'attimo}
  \item \emph{Say it:} \emph{l'attimo}, once more
  \item \emph{Say it:} a sentence of your own with \emph{l'attimo}
  \item \emph{Recall:} say the Italian for the length of time, then the Italian for the instant, then say \emph{l'attimo} again
\end{itemize}

\begin{tcolorbox}[breakable,colback=teal!4,colframe=teal!35!black,title={Before you move on}]
What does l'attimo mean? (\textquotedblleft{}The moment, a split second\textquotedblright{}.) Say it once more.
\end{tcolorbox}
\section[la fase]{la fase — the phase, the stage}

\label{lesson:IT-C277-fase}



\begin{quote}
Before the new one: say the Italian for the instant, then the Italian for the moment.
\end{quote}

\subsection*{You'll want to know: la fase}
\textbf{la fase} — \textquotedblleft{}the phase, the stage\textquotedblright{}.

\begin{quote}
Siamo nella prima fase del lavoro.
\end{quote}

\begin{quote}
\emph{We are in the first phase of the work.}
\end{quote}

Also the phases of the moon: \emph{le fasi della luna}.

\begin{cousinweb}
\emph{fase} is a learned borrowing from Greek \emph{phásis} \textquotedblleft{}appearance\textquotedblright{} (EtymDB), first said of the changing look of the moon.
\end{cousinweb}

\begin{grammarlens}[title={What you've built}]
One more word for telling, explaining and describing at B1.
\end{grammarlens}

\subsection*{Your turn}
\begin{itemize}
\raggedright
  \item \emph{Say it:} \emph{la fase}
  \item \emph{Say it:} \emph{la fase}, once more
  \item \emph{Say it:} a sentence of your own with \emph{la fase}
  \item \emph{Recall:} say the Italian for the instant, then the Italian for the moment, then say \emph{la fase} again
\end{itemize}

\begin{tcolorbox}[breakable,colback=teal!4,colframe=teal!35!black,title={Before you move on}]
What does la fase mean? (\textquotedblleft{}The phase, the stage\textquotedblright{}.) Say it once more.
\end{tcolorbox}
